% Section 5.1
\section{Análisis y Discusión}\label{sec:analisis-discusion}

Esta sección ofrece un análisis e interpretación de los resultados obtenidos en toda la investigación a partir de los estudios primarios seleccionados(SPSs), con el proposito de
dar respuesta a las preguntas de investigación(RQs) definidas en la fase de planeación. A diferencia de la sección de resultados, aquí no solo se describen los datos, sino que se examinan 
los patrones, tendencias y relaciones identificados entre los estudios análizados.

En terminos generales, se observa una distribución equilibrada de los resultados de los SPSs en relación con los RQs, lo que sugiere que la literatura aborda
ambas perspectivas de manera comparable. Este comportamiento evidencia un cierto grado de madurez y avance en el campo de investigación, en el cual se desarrollan tanto los fundamentos conceptuales,
reflejados en las metodologías identificadas, como los enfoques aplicados, evidenciados en las tecnologías implementadas en los distintos estudios.

Adicionalmente, el análisis muestra un incremento significativo en el número de publicaciones relacionados con la seguridad en entornos de seguridad en la nube durante los últimos años, particularmente 
alrededor de 2022 a 2025. Este crecimiento refleja un alto interés en este ambito investigativo, impulsado por la rapida evolución de las tecnologías y su adopción en los diferentes
sectores tanto empresariales como institucionales. En este contexto, la computación en la nube y la cibersegiridad se consolidan como elementos clave en el desarrollo y la protección de las infraestructuras digitales modernas.

\subsection{Discussion of RQ1}

En relación con la RQ1, los resultados evidencias que la mayoria de los estudios se centran en las áreas de seguridad, cibereseguridad, computación en la nube y controles de seguridad, esto indica que estos temas constituyen
el nucleo principal de la investigación. Esta predominancia sugiere que la literatura esta fuertemente orientada a la protección de información en infraestructuras digitales y en la gestión de riesgos en entornos de computación en la nube.
Asi mismo, se observa que los estudios no solo abordan aspectos teóricos, sino que tambien se proponen soluciones prácticas mediante el uso de diferentes herramientas tecnológicas y enfoques arquitectónicos. En este sentido, 
la presencia de herraminetas como Zero Trust, SIEM, Snort, IoT, inteligencia artificial y estándares como ISO/IEC 27001 muestra una tendencia hacia la implementación de marcos de seguridad robustos y adaptativos.
Por otro lado, aunque existe una amplia cobertura de estos temas dentro del estudio, aun existen zonas sin exploración abundante, creando vacios dentro de la investigación, que abre oportunidades a estudios e investigaciones que 
permitan fortalecer esas áreas dentro del dominio de la seguridad en la nube. Sin embargo, como se cuenta con una escasa integración de herramientas tecnológicas se sugiere como una necesidad a futuro el fortalecer la investigación aplicada, 
promoviendo el desarrollo de soluciones implementables que permitan validar y consolidar los enfoques teóricos existentes. En este sentido, se identifica una oportunidad clara para futuros trabajos orientados a cerrar la brecha entre la propuesta 
conceptual y su aplicación práctica, mediante el uso de herramientas que faciliten la implementación, evaluación y validación de los modelos planteados en la literatura.

\subsection{Discussion of RQ2}

En relación con la RQ2, los resultados muestran una clara predominancia de enfoques metodológicos sobre el uso de herramientas y si bien la cantidad de herramientas identificadas es reducida, se evidencia una recurrencia significativa de algunas de ellas 
a lo largo de diferentes estudios. Esto indica que existe una preferencia marcada por un conjunto específico de herramientas, las cuales pueden considerarse como las más representativas y consolidadas en el área.
La predominancia de los enfoques metodológicos se evidencia en la alta frecuencia de estudios centrados en propuestas, modelos, marcos conceptuales y taxonomías, este comportamiento sugiere que la investigación cuenta con una gran cantidad de aportes enfocados en 
el desarrollo de funfdamentos teóricos y estructurales. Adicionalmente, la presencia de estudios relacionados con clasificación, arquitectura, y taxonomia refuerza la idea de que aun se continua en proceso de organización y estructuración del conocimiento.
